\documentclass[11pt, a4paper]{report}

% ========== Common Packages ==========
\usepackage{fontspec}
\usepackage{kotex}
\usepackage{amsmath, amssymb, amsthm}
\usepackage{mathtools}
\usepackage{bm}
\usepackage{graphicx}
\usepackage{booktabs}
\usepackage{array}
\usepackage{tabularx}
\usepackage{longtable}
\usepackage{hyperref}
\usepackage{cleveref}
\usepackage{geometry}
\usepackage{fancyhdr}
\usepackage{titlesec}
\usepackage{enumitem}
\usepackage{xcolor}
\usepackage{tcolorbox}
\usepackage{algorithm}
\usepackage{algorithmic}
\usepackage{tikz}
\usetikzlibrary{arrows.meta, positioning, calc, decorations.pathreplacing, shapes.geometric, fit, patterns, angles, quotes, trees}
\usepackage{pgfplots}
\pgfplotsset{compat=1.18}
\usepackage{caption}
\usepackage{subcaption}
\usepackage{float}

\geometry{margin=2.5cm, top=3cm, bottom=3cm}

% ========== Colors ==========
% Common
\definecolor{defblue}{RGB}{30, 80, 150}
\definecolor{thmgreen}{RGB}{20, 110, 60}
\definecolor{noteorange}{RGB}{200, 100, 0}
\definecolor{lightblue}{RGB}{230, 240, 255}
\definecolor{lightgreen}{RGB}{230, 255, 235}
\definecolor{lightorange}{RGB}{255, 245, 230}
\definecolor{lightgray}{RGB}{245, 245, 245}
% Extended
\definecolor{lightred}{RGB}{255, 235, 235}
\definecolor{darkred}{RGB}{180, 30, 30}
\definecolor{biogreen}{RGB}{10, 130, 80}
\definecolor{cvpurple}{RGB}{100, 40, 150}
\definecolor{injuryred}{RGB}{200, 40, 40}
\definecolor{codegray}{RGB}{240, 240, 245}
\definecolor{pipeblue}{RGB}{25, 100, 180}
\definecolor{constraintred}{RGB}{200, 50, 50}
\definecolor{termblack}{RGB}{30, 30, 30}
\definecolor{goldcolor}{RGB}{180, 140, 20}

% ========== Theorem-like environments ==========
\theoremstyle{definition}
\newtheorem{definition}{Definition}[chapter]
\newtheorem{example}{Example}[chapter]

\theoremstyle{plain}
\newtheorem{theorem}{Theorem}[chapter]
\newtheorem{proposition}{Proposition}[chapter]

\theoremstyle{remark}
\newtheorem{remark}{Remark}[chapter]

% ========== tcolorbox environments ==========
% --- Common (all parts) ---
\newtcolorbox{keyidea}[1][]{
  colback=lightblue,
  colframe=defblue,
  fonttitle=\bfseries,
  title={Key Idea},
  #1
}

\newtcolorbox{importantnote}[1][]{
  colback=lightorange,
  colframe=noteorange,
  fonttitle=\bfseries,
  title={Important},
  #1
}

\newtcolorbox{summary}[1][]{
  colback=lightgreen,
  colframe=thmgreen,
  fonttitle=\bfseries,
  title={Summary},
  #1
}

% --- Warning/Error ---
\newtcolorbox{warningbox}[1][]{
  colback=lightred,
  colframe=darkred,
  fonttitle=\bfseries,
  title={Warning},
  #1
}

% --- Domain: Biomechanics & Clinical ---
\newtcolorbox{clinicalnote}[1][]{
  colback=lightred,
  colframe=darkred,
  fonttitle=\bfseries,
  title={Clinical Note},
  #1
}

\newtcolorbox{biomechbox}[1][]{
  colback=lightgreen,
  colframe=biogreen,
  fonttitle=\bfseries,
  title={Biomechanics},
  #1
}

\newtcolorbox{injurybox}[1][]{
  colback={injuryred!8},
  colframe=injuryred,
  fonttitle=\bfseries,
  title={Injury Risk},
  #1
}

% --- Domain: Computer Vision ---
\newtcolorbox{cvbox}[1][]{
  colback={cvpurple!5},
  colframe=cvpurple,
  fonttitle=\bfseries,
  title={Computer Vision},
  #1
}

% --- Pipeline & Setup ---
\newtcolorbox{setupbox}[1][]{
  colback=lightgray,
  colframe=gray!60,
  fonttitle=\bfseries,
  title={Setup},
  #1
}

\newtcolorbox{pipelinebox}[1][]{
  colback=pipeblue!5,
  colframe=pipeblue,
  fonttitle=\bfseries,
  title={Pipeline Step},
  #1
}

\newtcolorbox{codebox}[1][]{
  colback=codegray,
  colframe=gray!70,
  fonttitle=\bfseries\ttfamily,
  title={Code},
  #1
}

\newtcolorbox{termbox}[1][]{
  colback=termblack!5,
  colframe=termblack!60,
  fonttitle=\bfseries\ttfamily,
  title={Terminal},
  #1
}

% --- Research ---
\newtcolorbox{hypothesis}[1][]{
  colback=goldcolor!8,
  colframe=goldcolor,
  fonttitle=\bfseries,
  title={Hypothesis},
  #1
}

\newtcolorbox{resultbox}[1][]{
  colback=thmgreen!8,
  colframe=thmgreen,
  fonttitle=\bfseries,
  title={Expected Result},
  #1
}

% --- Constraints & Solutions ---
\newtcolorbox{constraintbox}[1][]{
  colback=constraintred!8,
  colframe=constraintred,
  fonttitle=\bfseries,
  title={Constraint},
  #1
}

\newtcolorbox{solutionbox}[1][]{
  colback=thmgreen!8,
  colframe=thmgreen,
  fonttitle=\bfseries,
  title={Solution},
  #1
}

\newtcolorbox{databox}[1][]{
  colback=pipeblue!5,
  colframe=pipeblue,
  fonttitle=\bfseries,
  title={Data Spec},
  #1
}

% --- Progress/Status ---
\newtcolorbox{successbox}[1][]{
  colback=lightgreen,
  colframe=thmgreen,
  fonttitle=\bfseries,
  title={Success},
  #1
}

\newtcolorbox{nextbox}[1][]{
  colback=lightorange,
  colframe=noteorange,
  fonttitle=\bfseries,
  title={Next Step},
  #1
}

% --- Fitting guide ---
\newtcolorbox{failbox}[1][]{
  colback=lightred,
  colframe=darkred,
  fonttitle=\bfseries,
  title={Failure Pattern},
  #1
}

\newtcolorbox{fixbox}[1][]{
  colback=lightgreen,
  colframe=thmgreen,
  fonttitle=\bfseries,
  title={Fix},
  #1
}

\newtcolorbox{lessonbox}[1][]{
  colback=lightorange,
  colframe=noteorange,
  fonttitle=\bfseries,
  title={Lesson Learned},
  #1
}

% --- Specification ---
\newtcolorbox{specbox}[1][]{
  colback=cvpurple!5,
  colframe=cvpurple,
  fonttitle=\bfseries,
  title={Specification},
  #1
}

\newtcolorbox{checklistbox}[1][]{
  colback=lightgreen,
  colframe=biogreen,
  fonttitle=\bfseries,
  title={Checklist},
  #1
}

% ========== Custom math commands ==========
% Vectors (lowercase)
\newcommand{\vx}{\mathbf{x}}
\newcommand{\vd}{\mathbf{d}}
\newcommand{\vo}{\mathbf{o}}
\newcommand{\vc}{\mathbf{c}}
\newcommand{\vr}{\mathbf{r}}
\newcommand{\vp}{\mathbf{p}}
\newcommand{\vv}{\mathbf{v}}
\newcommand{\vn}{\mathbf{n}}
\newcommand{\vu}{\mathbf{u}}
\newcommand{\vt}{\mathbf{t}}
\newcommand{\ve}{\mathbf{e}}
\newcommand{\vq}{\mathbf{q}}
% Vectors (uppercase)
\newcommand{\vF}{\mathbf{F}}
\newcommand{\vM}{\mathbf{M}}
\newcommand{\va}{\mathbf{a}}
\newcommand{\vg}{\mathbf{g}}
\newcommand{\vX}{\mathbf{X}}
% Bold Greek
\newcommand{\vmu}{\bm{\mu}}
\newcommand{\vbeta}{\bm{\beta}}
\newcommand{\vtheta}{\bm{\theta}}
\newcommand{\vpsi}{\bm{\psi}}
% Matrices
\newcommand{\mSigma}{\bm{\Sigma}}
\newcommand{\mR}{\mathbf{R}}
\newcommand{\mS}{\mathbf{S}}
\newcommand{\mW}{\mathbf{W}}
\newcommand{\mJ}{\mathbf{J}}
\newcommand{\mG}{\mathbf{G}}
\newcommand{\mT}{\mathbf{T}}
\newcommand{\mI}{\mathbf{I}}
\newcommand{\mB}{\mathbf{B}}
\newcommand{\mK}{\mathbf{K}}
\newcommand{\mP}{\mathbf{P}}
\newcommand{\mF}{\mathbf{F}}
\newcommand{\mE}{\mathbf{E}}
\newcommand{\mA}{\mathbf{A}}
\newcommand{\mH}{\mathbf{H}}
% Sets and groups
\newcommand{\RR}{\mathbb{R}}
\newcommand{\SO}{\mathrm{SO}}
% Units
\newcommand{\degps}{\,\text{deg/s}}
\newcommand{\Nm}{\ensuremath{\,\text{N}\cdot\text{m}}}
\newcommand{\BW}{\,\text{BW}}

% ========== Listings ==========
\usepackage{listings}

\lstset{
  basicstyle=\ttfamily\small,
  keywordstyle=\color{defblue}\bfseries,
  commentstyle=\color{thmgreen}\itshape,
  stringstyle=\color{noteorange},
  backgroundcolor=\color{codegray},
  frame=single,
  rulecolor=\color{gray!40},
  breaklines=true,
  showstringspaces=false,
  numbers=left,
  numberstyle=\tiny\color{gray},
  tabsize=4,
  captionpos=b
}

% ========== Header/Footer ==========
\pagestyle{fancy}
\fancyhf{}
\fancyhead[L]{\leftmark}
\fancyhead[R]{\thepage}
\renewcommand{\headrulewidth}{0.4pt}

% ========== Title formatting ==========
\titleformat{\chapter}[display]
  {\normalfont\huge\bfseries\color{defblue}}
  {\chaptertitlename\ \thechapter}{20pt}{\Huge}

\titleformat{\section}
  {\normalfont\Large\bfseries\color{defblue!80!black}}
  {\thesection}{1em}{}

\titleformat{\subsection}
  {\normalfont\large\bfseries\color{defblue!60!black}}
  {\thesubsection}{1em}{}


\begin{document}

% ==================== TITLE PAGE ====================
\begin{titlepage}
\centering
\vspace*{2cm}
{\Huge\bfseries\color{defblue} SAM 3D Body 기반\\[0.3cm] 투구 분석 파이프라인}

\vspace{1cm}
{\LARGE\color{defblue!70!black} 설치 $\to$ 추론 $\to$ 변환 $\to$ GART $\to$ OpenSim}

\vspace{2cm}
{\large 7-Camera 240Hz 마커리스 투구 생체역학 분석}

\vspace{1cm}
{\large 2026년 4월 14일}

\vspace{3cm}
\begin{tcolorbox}[colback=lightblue, colframe=defblue, width=0.85\textwidth]
\textbf{핵심}: SAM 3D Body (PA-MPJPE 30.4mm, SOTA)로 7개 카메라 각각 추론\\
$\to$ MHR 메시 $\to$ SMPL 변환 $\to$ 7-view reproj 정제 $\to$ GART 포토메트릭 정제\\
$\to$ OpenSim 생체역학 분석. \textbf{총 소요: 2--3시간.}
\end{tcolorbox}
\end{titlepage}

\tableofcontents
\newpage


% ====================================================================
\chapter{왜 SAM 3D Body인가}
% ====================================================================

\section{현재 파이프라인의 병목}

\begin{center}
\renewcommand{\arraystretch}{1.3}
\begin{tabular}{lccl}
\toprule
\textbf{문제} & \textbf{현재값} & \textbf{목표} & \textbf{원인} \\
\midrule
Reproj 오차 & 153.5px & $<$50px & 캘리브레이션 + 정규화 \\
GART 학습 & 3K 스텝 & 30K & 부족한 반복 \\
생체역학 & 미구현 & 관절각 $<$5° & 코드 공백 \\
\bottomrule
\end{tabular}
\end{center}

\begin{warningbox}
\textbf{냉정한 진단}: 153px 오차의 바닥(113px)이 body\_pose 정규화 때문인지 캘리브레이션 부정확 때문인지 구분이 안 됩니다. VGGT 캘리브레이션은 체커보드 없이 AI 추정이며, COLMAP은 카메라 간 초점거리 편차가 2.6배입니다.
\end{warningbox}

\section{SAM 3D Body의 이점}

\begin{center}
\renewcommand{\arraystretch}{1.3}
\begin{tabular}{lcc}
\toprule
& \textbf{U-HMR (현재)} & \textbf{SAM 3D Body} \\
\midrule
입력 & 4-view 동시 필수 & \textbf{단안 (카메라별 독립)} \\
출력 모델 & SMPL (6,890 vtx) & \textbf{MHR (18,439 vtx)} \\
정확도 (3DPW) & 31.0mm MPJPE & \textbf{30.4mm PA-MPJPE} \\
아키텍처 제약 & 정확히 4-view & \textbf{제한 없음 (1--N view)} \\
캘리브레이션 의존 & 다시점 융합 필요 & \textbf{불필요 (단안)} \\
카메라 호환 & 동일 해상도만 & \textbf{임의 해상도 (cam1 세로 포함)} \\
라이선스 & SMPL (비상업) & \textbf{MHR Apache 2.0} \\
설치 시간 & 이미 완료 & \textbf{20--30분} \\
\bottomrule
\end{tabular}
\end{center}

\begin{keyidea}
\textbf{핵심 이점}: SAM 3D Body는 \textbf{캘리브레이션 없이} 각 카메라 독립 추론이 가능합니다. 현재 파이프라인의 최대 병목인 캘리브레이션 품질 문제를 \textbf{우회}합니다. 7대 카메라 각각에서 MHR 메시를 독립 추론한 뒤, 다시점 정합은 GART 포토메트릭 손실에 맡깁니다.
\end{keyidea}


% ====================================================================
\chapter{다시점 3DGS가 투구 정확도를 높이는 원리}
% ====================================================================

3DGS(GART)는 단순히 ``예쁜 렌더링''이 아닙니다. \textbf{모션 품질 자체를 근본적으로 개선}하는 메커니즘입니다. 이 장에서는 왜 다시점 3DGS가 투구 분석에서 기존 키포인트 기반 방법을 넘어서는지 분석합니다.

\section{감독 신호 밀도: 124,000배 차이}

기존 키포인트 기반 감독과 3DGS 포토메트릭 감독의 정보량을 비교합니다.

\begin{center}
\renewcommand{\arraystretch}{1.3}
\begin{tabular}{lccc}
\toprule
\textbf{방법} & \textbf{값/프레임} & \textbf{대 SMPL 85파라미터} & \textbf{과결정 비율} \\
\midrule
단안 키포인트 (25$\times$2) & 50 & 50:85 & 0.6:1 (\textbf{부족}) \\
7-view 키포인트 (25$\times$2$\times$7) & 350 & 350:85 & 4.1:1 \\
단안 3DGS (1920$\times$1080$\times$3) & 6.2M & 6.2M:85 & 73,000:1 \\
\textbf{7-view 3DGS} & \textbf{43.5M} & \textbf{43.5M:85} & \textbf{512,000:1} \\
\bottomrule
\end{tabular}
\end{center}

\begin{keyidea}
키포인트 기반: 25개 관절의 \textbf{위치}만 감독합니다 (점).\\
포토메트릭 기반: 인체의 \textbf{전체 표면}을 감독합니다 (면).\\
관절 사이의 팔뚝 형태, 유니폼 주름, 근육 윤곽, 실루엣 경계가 모두 감독 신호입니다.
\end{keyidea}


\section{메커니즘 1: 다시점 기하학적 제약}

키포인트 기반에서는 포즈 $\vtheta$가 한 카메라의 2D 투영에만 맞으면 됩니다. 깊이 축 오류는 처벌되지 않습니다.

\textbf{다시점 3DGS에서는 포즈 $\vtheta$로 렌더링한 이미지가 7개 카메라 \underline{모두}에서 동시에 맞아야 합니다.}

\begin{equation}
\vtheta^* = \arg\min_{\vtheta} \sum_{c=1}^{7} \Big[ (1-\lambda_s)\|I_c - \hat{I}_c(\vtheta)\|_1 + \lambda_s \mathcal{L}_{\text{SSIM}}(I_c, \hat{I}_c(\vtheta)) \Big]
\end{equation}

잘못된 포즈는 \textbf{어떤 각도에서 반드시 틀리게 보입니다}:

\begin{center}
\begin{tikzpicture}[
  node distance=1.5cm,
  cam/.style={draw, fill=lightblue, minimum width=1.5cm, minimum height=0.6cm, align=center, font=\scriptsize},
  body/.style={draw, fill=lightorange, circle, minimum size=1cm, align=center, font=\scriptsize},
  arrow/.style={->, thick, >=stealth, dashed}
]
\node[body] (b) {투구자\\포즈 $\vtheta$};
\node[cam, above left=1.5cm and 2cm of b] (c1) {cam1\\(후면)};
\node[cam, above=1.8cm of b] (c3) {cam3\\(정면)};
\node[cam, above right=1.5cm and 2cm of b] (c5) {cam5\\(측면)};
\node[cam, left=2.5cm of b] (c2) {cam2\\(좌45°)};
\node[cam, right=2.5cm of b] (c7) {cam7\\(우측면)};

\draw[arrow] (b) -- (c1) node[midway, above, font=\tiny] {팔 보임};
\draw[arrow] (b) -- (c3) node[midway, right, font=\tiny] {정면};
\draw[arrow] (b) -- (c5) node[midway, above, font=\tiny] {깊이 확인};
\draw[arrow] (b) -- (c2) node[midway, above, font=\tiny] {45° 검증};
\draw[arrow] (b) -- (c7) node[midway, above, font=\tiny] {팔꿈치 각도};
\end{tikzpicture}
\end{center}

\begin{importantnote}
\textbf{예시}: 어깨 외회전(MER) 170°와 150°는 정면(cam3)에서 키포인트가 거의 동일합니다. 하지만 측면(cam5, cam7)에서는 상완의 \textbf{실루엣}이 완전히 달라집니다. 다시점 3DGS는 이 차이를 수백만 픽셀로 감지합니다.
\end{importantnote}


\section{메커니즘 2: 깊이 복원 --- 포토메트릭 깊이 제약}

\subsection{키포인트 삼각측량의 한계}

키포인트 기반 깊이 복원(DLT 삼각측량):
\begin{equation}
\sigma_z \propto \frac{z^2}{f \cdot B}
\end{equation}
여기서 $B$는 기선(baseline) 길이. 좁은 기선 $\to$ 깊이 불확실성 폭발. 또한 \textbf{키포인트 노이즈 5--15px가 깊이 오차로 직접 전파}됩니다.

\subsection{3DGS 포토메트릭 깊이 제약}

3DGS에서는 깊이 오류가 \textbf{실루엣 불일치}로 드러납니다:

\begin{center}
\renewcommand{\arraystretch}{1.3}
\begin{tabular}{lcc}
\toprule
\textbf{움직임} & \textbf{정면(cam3) 변화} & \textbf{측면(cam5) 변화} \\
\midrule
팔 z축 10cm 전진 & 미미 (0--2px) & \textbf{$\sim$50px 이동} \\
어깨 내회전 20° & 미미 (키포인트 불변) & \textbf{실루엣 형태 변화} \\
팔꿈치 외반 5° & 미미 & \textbf{45° 뷰에서 감지} \\
\bottomrule
\end{tabular}
\end{center}

\textbf{각 카메라가 다른 축의 오류에 민감} $\to$ 7개 조합하면 모든 방향의 오류가 감지됩니다.


\section{메커니즘 3: 자기 폐색(self-occlusion) 해소}

투구의 6단계 중 여러 순간에서 몸통이 투구 팔을 가립니다:

\begin{center}
\renewcommand{\arraystretch}{1.3}
\begin{tabular}{lll}
\toprule
\textbf{단계} & \textbf{폐색 상황} & \textbf{정면에서 보이는 것} \\
\midrule
와인드업 & 글러브가 팔 앞에 위치 & 팔 안 보임 \\
레이트 코킹 & 팔이 몸 뒤로 $\to$ 정면 폐색 & 어깨만 보임 \\
가속기 & 팔이 머리 뒤에서 앞으로 & 순간적 폐색 \\
\bottomrule
\end{tabular}
\end{center}

\paragraph{키포인트 기반 (단안/다시점).}
폐색된 뷰에서 키포인트 신뢰도 = 0 $\to$ 해당 관절의 정보 \textbf{완전 소실}.

\paragraph{다시점 3DGS.}
\begin{itemize}[leftmargin=*]
  \item 폐색된 뷰(cam3): ``팔이 안 보이는 것'' 자체가 정답 렌더링 $\to$ \textbf{실루엣 경계가 감독 신호}
  \item 비폐색 뷰(cam7 후면): 팔이 보임 $\to$ \textbf{실제 팔 자세를 감독}
  \item 두 정보를 결합하면 폐색 순간에도 정확한 포즈 추정 가능
\end{itemize}

\begin{keyidea}
다시점 3DGS에서는 \textbf{보이지 않는 것도 정보}입니다. ``이 각도에서 팔이 보이면 안 됨''이라는 제약이 포즈를 강하게 결정합니다.
\end{keyidea}


\section{메커니즘 4: 모션 블러가 정보가 됨}

투구 가속기에서 7,000°/s 어깨 회전 $\to$ 1/1000s 셔터에서 프레임당 7° 블러.

\paragraph{키포인트 기반.}
블러 $\to$ 키포인트 검출기 혼란 $\to$ 신뢰도 급감 $\to$ \textbf{가장 중요한 순간에 가장 부정확}.

\paragraph{3DGS 포토메트릭.}
블러된 이미지 자체를 재현하도록 가우시안을 학습시킵니다. 블러 방향은 움직임 방향과 일치하므로:
\begin{itemize}[leftmargin=*]
  \item 블러 길이 $\propto$ 각속도 $\to$ 속도 정보
  \item 블러 방향 = 회전축에 수직 $\to$ 회전축 정보
  \item 여러 뷰의 블러 패턴 조합 $\to$ 3D 회전 벡터 복원 가능
\end{itemize}

\begin{importantnote}
이것이 GART-Pitch의 \textbf{핵심 기여점}입니다. 기존 모든 마커리스 투구 시스템(Theia3D 52mm, KinaTrax 57mm)은 키포인트 기반입니다. 포토메트릭 3DGS 포즈 정제를 투구에 적용한 연구는 \textbf{세계적으로 0건}입니다.
\end{importantnote}


\section{투구 위상별 3DGS 기여도}

\begin{center}
\renewcommand{\arraystretch}{1.3}
\begin{tabular}{lccc}
\toprule
\textbf{투구 위상} & \textbf{키포인트 품질} & \textbf{3DGS 기여도} & $\lambda_{\text{temp}}$ \\
\midrule
와인드업 (느림) & 양호 & 보완적 & 10.0 \\
스트라이드 & 양호 & 중간 & 5.0 \\
레이트 코킹 & 저하 (폐색 시작) & \textbf{높음} & 2.0 \\
\textbf{가속기 (30ms)} & \textbf{매우 나쁨 (블러)} & \textbf{결정적} & \textbf{0.5} \\
감속 & 중간 & 높음 & 1.0 \\
팔로쓰루 & 양호 & 보완적 & 3.0 \\
\bottomrule
\end{tabular}
\end{center}

\begin{summary}
\textbf{다시점 3DGS가 투구 정확도를 높이는 4가지 메커니즘}:
\begin{enumerate}
\item \textbf{기하학적 제약}: 7개 뷰에서 동시에 맞아야 함 $\to$ 포즈가 유일하게 결정됨 (과결정 512,000:1)
\item \textbf{깊이 복원}: 각 카메라가 다른 축 오류에 민감 $\to$ 실루엣으로 깊이 제약
\item \textbf{폐색 해소}: 보이지 않는 것도 정보 $\to$ 비폐색 뷰와 결합하여 폐색 순간 포즈 추정
\item \textbf{블러 활용}: 블러 패턴이 속도/방향 정보 $\to$ 가속기에서 키포인트가 실패하는 순간에 3DGS가 주도
\end{enumerate}
SAM 3D Body(좋은 초기값) + 다시점 3DGS(정밀 정제) = 기존 SOTA를 넘을 수 있는 구체적 메커니즘.
\end{summary}


\section{SAM 3D Body + 다시점 3DGS 시너지}

\begin{center}
\begin{tikzpicture}[
  node distance=1.0cm,
  box/.style={draw, rounded corners, minimum width=5cm, minimum height=0.8cm, align=center, font=\small},
  arrow/.style={->, thick, >=stealth}
]
\node[box, fill=lightorange] (sam) {SAM 3D Body (단안 SOTA 30.4mm)\\7개 카메라 독립 추론};
\node[box, fill=lightblue, below=of sam] (init) {좋은 초기값 (7-view 중위수)\\캘리브레이션 의존 최소화};
\node[box, fill=lightgreen, below=of init] (gart) {GART 7-view 가우시안 학습\\30K스텝, 다시점 일관성};
\node[box, fill={cvpurple!15}, below=of gart] (refine) {포토메트릭 포즈 정제\\43.5M 픽셀 감독 $\times$ 위상 적응형};
\node[box, fill=lightblue, below=of refine] (result) {최종 SMPL 포즈\\MPJPE $<$20mm 목표 (SOTA 52mm 대비 60\%↓)};

\draw[arrow] (sam) -- (init) node[midway, right, font=\scriptsize] {MHR→SMPL 변환};
\draw[arrow] (init) -- (gart) node[midway, right, font=\scriptsize] {수렴 가속};
\draw[arrow] (gart) -- (refine) node[midway, right, font=\scriptsize] {가우시안 고정, $\vtheta$만 최적화};
\draw[arrow] (refine) -- (result) node[midway, right, font=\scriptsize] {가속기에서 $\lambda_{\text{temp}}$=0.5};
\end{tikzpicture}
\end{center}

\begin{keyidea}
\textbf{시너지}: SAM 3D Body가 좋은 초기값 제공 $\to$ GART 학습이 빨리 수렴 $\to$ 포즈 정제가 더 정확한 영역에서 시작 $\to$ \textbf{최종 포즈 품질 최대화}. 이것이 ``3DGS를 쓴다''가 아니라 ``SAM 3D Body + 다시점 3DGS의 조합이 SOTA를 넘는 구체적 경로''입니다.
\end{keyidea}


% ====================================================================
\chapter{전체 파이프라인}
% ====================================================================

\section{아키텍처 개요}

\begin{center}
\begin{tikzpicture}[
  node distance=1.0cm,
  box/.style={draw, rounded corners, fill=lightblue, minimum width=5.5cm, minimum height=0.8cm, align=center, font=\small},
  newbox/.style={draw, rounded corners, fill=lightorange, minimum width=5.5cm, minimum height=0.8cm, align=center, font=\small\bfseries},
  existbox/.style={draw, rounded corners, fill=lightgreen, minimum width=5.5cm, minimum height=0.8cm, align=center, font=\small},
  arrow/.style={->, thick, >=stealth},
  label/.style={font=\scriptsize, midway, right, text width=4cm}
]

\node[box] (input) {7-cam $\times$ 1000프레임 (240Hz)};
\node[newbox, below=of input] (sam3d) {SAM 3D Body $\times$ 7\\각 카메라 독립 MHR 추론};
\node[newbox, below=of sam3d] (convert) {MHR $\to$ SMPL 변환\\pymomentum / PyTorch 최적화};
\node[existbox, below=of convert] (reproj) {7-View Reproj Fitting\\$\lambda=0.0001$, 500스텝};
\node[existbox, below=of reproj] (gart) {GART 7-View 30K스텝\\+ 포토메트릭 포즈 정제};
\node[newbox, below=of gart] (bio) {SMPL $\to$ TRC $\to$ OpenSim IK\\관절 각도 + 토크};
\node[box, below=of bio] (valid) {Vicon 검증 (동기화 완료)\\MPJPE, Bland-Altman};

\draw[arrow] (input) -- (sam3d) node[label] {\textbf{NEW}: 설치 20분};
\draw[arrow] (sam3d) -- (convert) node[label] {\textbf{NEW}: MHR→SMPL};
\draw[arrow] (convert) -- (reproj) node[label] {기존 스크립트 재활용};
\draw[arrow] (reproj) -- (gart) node[label] {기존 스크립트 재활용};
\draw[arrow] (gart) -- (bio) node[label] {\textbf{NEW}: SAM4Dcap 마커 매핑};
\draw[arrow] (bio) -- (valid) node[label] {로컬에서 동기화 완료};
\end{tikzpicture}
\end{center}

\begin{importantnote}
\textcolor{noteorange}{\textbf{주황색}} = 새로 추가하는 단계 (3개)\\
\textcolor{thmgreen}{\textbf{초록색}} = 기존 스크립트 재활용 (2개)\\
기존 U-HMR $\to$ SAM 3D Body로 교체하면서 캘리브레이션 의존성을 제거합니다.
\end{importantnote}


\section{현재 vs 새 파이프라인 비교}

\begin{center}
\renewcommand{\arraystretch}{1.3}
\begin{tabular}{p{3cm}p{5cm}p{5cm}}
\toprule
\textbf{단계} & \textbf{현재} & \textbf{SAM 3D Body 기반} \\
\midrule
초기 포즈 & U-HMR 4-view 융합 & SAM 3D Body $\times$7 (독립) \\
바디 모델 & SMPL 직접 & MHR $\to$ SMPL 변환 \\
캘리브레이션 & VGGT (AI 추정) & \textbf{추론에 불필요} \\
cam1 (세로) & 제외 & \textbf{포함 (그대로 입력)} \\
Reproj 정제 & 4-view, $\lambda$=0.01 & 7-view, $\lambda$=0.0001 \\
GART & 4-view 3K 스텝 & 7-view 30K 스텝 \\
생체역학 & 미구현 & SAM4Dcap 마커 매핑 \\
검증 & Vicon 미동기화 & \textbf{동기화 완료 (125프레임)} \\
\bottomrule
\end{tabular}
\end{center}


% ====================================================================
\chapter{Step 1: SAM 3D Body 설치 (20분)}
% ====================================================================

\section{Conda 환경 생성}

\begin{codebox}[title={서버 명령어}]
\begin{verbatim}
# 새 환경 (GART와 충돌 방지)
conda create -n sam3d python=3.11 -y
conda activate sam3d

# PyTorch (CUDA 11.8 — 서버 기준)
pip install torch torchvision torchaudio \
    --index-url https://download.pytorch.org/whl/cu118

# 의존성 한 번에
pip install pytorch-lightning pyrender opencv-python yacs \
    scikit-image einops timm dill pandas rich \
    hydra-core pyrootutils webdataset roma joblib \
    seaborn loguru optree fvcore pycocotools \
    tensorboard huggingface_hub networkx==3.2.1

# Detectron2 (소스 빌드, 5-10분)
pip install 'git+https://github.com/facebookresearch/detectron2.git@a1ce2f9' \
    --no-build-isolation --no-deps
\end{verbatim}
\end{codebox}

\section{체크포인트 다운로드}

\begin{codebox}[title={HuggingFace 다운로드 (2.8GB, 1-5분)}]
\begin{verbatim}
# HuggingFace 로그인 (토큰 필요)
huggingface-cli login

# 게이트 승인: 브라우저에서 아래 URL 방문 → "Agree" 클릭
# https://huggingface.co/facebook/sam-3d-body-dinov3

# 다운로드
huggingface-cli download facebook/sam-3d-body-dinov3 \
    --local-dir ~/sam-3d-body/checkpoints/sam-3d-body-dinov3

# 확인
ls -la ~/sam-3d-body/checkpoints/sam-3d-body-dinov3/
# model.ckpt (2.11GB), assets/mhr_model.pt (696MB)
\end{verbatim}
\end{codebox}

\section{SAM 3D Body 코드 클론}

\begin{codebox}[title={레포 클론}]
\begin{verbatim}
cd ~
git clone https://github.com/facebookresearch/sam-3d-body.git
cd sam-3d-body
\end{verbatim}
\end{codebox}

\section{설치 확인 (1장 테스트)}

\begin{codebox}[title={단일 이미지 테스트}]
\begin{verbatim}
cd ~/sam-3d-body
export PYOPENGL_PLATFORM=egl  # 서버 headless 환경

python -c "
from notebook.utils import setup_sam_3d_body
import cv2
estimator = setup_sam_3d_body(
    hf_repo_id='facebook/sam-3d-body-dinov3',
    local_dir='checkpoints/sam-3d-body-dinov3'
)
img = cv2.imread('/home/elicer/images/cam3/00001.png')
out = estimator.process_one_image(cv2.cvtColor(img, cv2.COLOR_BGR2RGB))
print('Output keys:', list(out.keys()))
print('Vertices shape:', out['pred_vertices'].shape)  # (18439, 3)
print('SUCCESS')
"
\end{verbatim}
\end{codebox}

\begin{successbox}
설치 완료 기준: \texttt{Vertices shape: (18439, 3)} 출력 + \texttt{SUCCESS} 메시지.
\end{successbox}


% ====================================================================
\chapter{Step 2: 7-Camera 배치 추론 (20-40분)}
% ====================================================================

\section{추론 전략}

\begin{keyidea}
7개 카메라 $\times$ 999 프레임 = 6,993장을 SAM 3D Body로 독립 추론합니다.
각 이미지에서 MHR 파라미터 (18,439 꼭짓점 메시)를 추출합니다.
\textbf{캘리브레이션 불필요} --- 각 이미지가 독립적으로 처리됩니다.
\end{keyidea}

\section{배치 추론 스크립트}

\begin{codebox}[title={run\_sam3d\_7cam.py}]
\begin{verbatim}
"""7-camera x 999 frames SAM 3D Body batch inference."""
import os, cv2, numpy as np, torch, time
from notebook.utils import setup_sam_3d_body

estimator = setup_sam_3d_body(
    hf_repo_id='facebook/sam-3d-body-dinov3',
    local_dir='checkpoints/sam-3d-body-dinov3'
)

CAM_DIRS = {
    'cam1': '/home/elicer/images/cam1',   # 세로 1080x1920 → 그대로 입력!
    'cam2': '/home/elicer/images/cam2',
    'cam3': '/home/elicer/images/cam3',
    'cam4': '/home/elicer/images/cam4',
    'cam5': '/home/elicer/images/cam5',
    'cam6': '/home/elicer/images/cam6',
    'cam7': '/home/elicer/images/cam7',
}
OUT_DIR = '/home/elicer/sam3d_results'
N_FRAMES = 999

for cam_name, cam_dir in CAM_DIRS.items():
    out_cam = os.path.join(OUT_DIR, cam_name)
    os.makedirs(out_cam, exist_ok=True)

    files = sorted([f for f in os.listdir(cam_dir)
                    if f.endswith(('.jpg','.png'))])[:N_FRAMES]

    t0 = time.time()
    for fi, fn in enumerate(files):
        img = cv2.imread(os.path.join(cam_dir, fn))
        out = estimator.process_one_image(
            cv2.cvtColor(img, cv2.COLOR_BGR2RGB))

        np.savez_compressed(
            os.path.join(out_cam, f'{fi:06d}.npz'),
            pred_vertices=out['pred_vertices'],   # (18439, 3)
            focal_length=out.get('focal_length', [0]),
            pred_cam_t=out.get('pred_cam_t', [0,0,0]),
        )

        if fi % 100 == 0:
            elapsed = time.time() - t0
            fps = (fi+1) / elapsed if elapsed > 0 else 0
            eta = (len(files)-fi-1) / fps if fps > 0 else 0
            print(f'{cam_name} {fi}/{len(files)}: '
                  f'{fps:.1f} fps, ETA {eta:.0f}s')

    total = time.time() - t0
    print(f'{cam_name}: {len(files)} frames in {total:.0f}s '
          f'({len(files)/total:.1f} fps)')
\end{verbatim}
\end{codebox}

\section{예상 시간}

\begin{center}
\renewcommand{\arraystretch}{1.2}
\begin{tabular}{lccc}
\toprule
\textbf{추론 속도} & \textbf{프레임 수} & \textbf{시간} & \textbf{비고} \\
\midrule
$\sim$1 FPS (A100) & 6,993 & $\sim$117분 & 전체 순차 \\
$\sim$2 FPS (배치) & 6,993 & $\sim$58분 & 배치 최적화 \\
\midrule
\textbf{현실적 예상} & & \textbf{60--120분} & \\
\bottomrule
\end{tabular}
\end{center}

\begin{importantnote}
\textbf{시간 단축 전략}: 처음에는 \textbf{cam3 1개 $\times$ 100프레임}만 테스트 (2분).\\
품질 확인 후 전체 7-cam 실행. 또는 cam3,5,6,7 (4개)만 먼저 실행하여 비교.
\end{importantnote}


% ====================================================================
\chapter{Step 3: MHR $\to$ SMPL 변환 (10-20분)}
% ====================================================================

\section{변환 방법}

SAM 3D Body는 MHR 메시(18,439 꼭짓점)를 출력합니다. GART는 SMPL(6,890 꼭짓점)을 필요로 합니다. 변환이 필요합니다.

\subsection{방법 A: MHR 레포의 공식 변환 도구}

\begin{codebox}[title={MHR $\to$ SMPL 변환}]
\begin{verbatim}
# MHR 레포 클론
git clone https://github.com/facebookresearch/MHR.git
cd MHR

# pymomentum 설치 (CPU 충분)
pip install pymomentum-cpu

# 변환 실행
python tools/mhr_smpl_conversion/conversion.py \
    --input_dir /home/elicer/sam3d_results/cam3 \
    --output_dir /home/elicer/sam3d_smpl/cam3 \
    --direction mhr_to_smpl \
    --backend pytorch  # GPU 사용 (빠름)
\end{verbatim}
\end{codebox}

\begin{warningbox}
\textbf{함정 F-6 (cm/m 스케일)}: SAM 3D Body는 미터 단위로 출력하지만, MHR 변환 코드의 일부 버전은 센티미터를 가정합니다. \texttt{conversion.py}에서 \texttt{0.01} 스케일링 팩터가 있으면 제거하세요.
\end{warningbox}

\subsection{방법 B: 직접 꼭짓점 피팅 (대안)}

MHR 메시 꼭짓점 $\to$ SMPL 파라미터를 직접 최적화:

\begin{codebox}[title={MHR vertices $\to$ SMPL params 직접 피팅}]
\begin{verbatim}
import torch, smplx, numpy as np

model = smplx.create(model_path, model_type='smpl',
                     gender='neutral').cuda()

for fi in range(N_FRAMES):
    mhr = np.load(f'sam3d_results/cam3/{fi:06d}.npz')
    target_verts = torch.tensor(
        mhr['pred_vertices'], device='cuda', dtype=torch.float32)

    # SMPL 파라미터 최적화
    body_pose = torch.zeros(1, 69, requires_grad=True, device='cuda')
    betas = torch.zeros(1, 10, requires_grad=True, device='cuda')
    global_orient = torch.zeros(1, 3, requires_grad=True, device='cuda')
    transl = torch.zeros(1, 3, requires_grad=True, device='cuda')

    opt = torch.optim.Adam([body_pose, betas, global_orient, transl],
                           lr=0.01)

    for step in range(200):
        out = model(body_pose=body_pose, betas=betas,
                    global_orient=global_orient, transl=transl)
        # SMPL 6890 vertices → MHR 18439과 대응점 매칭 필요
        # 최근접점 매칭 또는 사전 계산된 대응 테이블 사용
        loss = ((out.vertices[0] - target_verts_sampled)**2).sum()
        opt.zero_grad(); loss.backward(); opt.step()
\end{verbatim}
\end{codebox}

\begin{importantnote}
\textbf{권장}: 방법 A (공식 도구)를 먼저 시도. 실패 시 방법 B로 전환.\\
공식 도구는 barycentric interpolation으로 topology 매핑을 정확히 처리합니다.
\end{importantnote}


% ====================================================================
\chapter{Step 4: 다시점 정합 + GART (기존 스크립트)}
% ====================================================================

\section{SAM 3D Body 출력 $\to$ Reproj Fitting}

SAM 3D Body는 각 카메라에서 독립적으로 SMPL body\_pose를 추정합니다. 7개 카메라의 결과를 통합해야 합니다.

\subsection{전략 A: 최선 뷰 선택 + Reproj (빠름)}

\begin{codebox}[title={7-view 평균 body\_pose $\to$ reproj 정제}]
\begin{verbatim}
# 7개 카메라의 body_pose를 평균
bp_views = []  # (7, 69) 각 카메라의 body_pose
for cam in ['cam1','cam2','cam3','cam4','cam5','cam6','cam7']:
    smpl = load_smpl_params(f'sam3d_smpl/{cam}/{fi:06d}.npz')
    bp_views.append(smpl['body_pose'])

# 중위수 (outlier에 강건)
body_pose_init = np.median(np.array(bp_views), axis=0)

# 이 초기값으로 7-view reproj fitting 실행
# → 02_reproj_7view_sequential.py 의 body_pose_init 교체
\end{verbatim}
\end{codebox}

\subsection{전략 B: 직접 GART 입력 (캘리브레이션 우회)}

\begin{keyidea}
SAM 3D Body가 \textbf{각 카메라의 카메라 좌표계에서} SMPL을 추정하므로, 캘리브레이션 없이 각 뷰의 SMPL을 직접 GART에 입력할 수 있습니다. GART는 각 (프레임, 카메라) 쌍에 대해 독립적으로 렌더링하므로, 월드 좌표계 통합이 불필요합니다.
\end{keyidea}

\begin{warningbox}
단, 이 경우 GART가 \textbf{캐노니컬 공간에서 7개 뷰를 일관되게} 학습해야 하므로, body\_pose는 동일하고 Rh/Th만 카메라별로 다르게 설정해야 합니다. 이를 위해서는 최소한의 상대 포즈(relative pose) 정보가 필요합니다.
\end{warningbox}


\section{GART 7-View 학습}

Step 3에서 얻은 SMPL 파라미터로 GART를 학습합니다. 기존 스크립트 (\texttt{03\_prepare\_gart\_data.py})를 그대로 사용합니다.

\begin{center}
\renewcommand{\arraystretch}{1.2}
\begin{tabular}{lcc}
\toprule
\textbf{설정} & \textbf{현재} & \textbf{SAM 3D Body 기반} \\
\midrule
초기 SMPL & U-HMR 4-view & \textbf{SAM 3D Body 7-view 중위수} \\
카메라 수 & 4 & \textbf{7} \\
학습 스텝 & 3,000 & \textbf{30,000} \\
포즈 정제 시작 & 미적용 & \textbf{스텝 15,000} \\
IMAGE\_ZOOM & 0.5 & \textbf{0.25} \\
\bottomrule
\end{tabular}
\end{center}


% ====================================================================
\chapter{Step 5: 생체역학 분석 (SAM4Dcap 활용)}
% ====================================================================

\section{SMPL $\to$ 해부학적 마커}

SAM4Dcap에서 추출한 SMPL 꼭짓점 $\to$ 마커 매핑 테이블을 사용합니다.

\begin{codebox}[title={핵심 마커 인덱스 (SMPL 6,890 꼭짓점 기준)}]
\begin{verbatim}
# SAM4Dcap updated_mapping.csv에서 추출
MARKER_MAP = {
    'C7':    1306,   # 경추 7번
    'LSHO':  1239,   # 왼쪽 어깨
    'RSHO':  4878,   # 오른쪽 어깨
    'RELB':  5090,   # 오른쪽 팔꿈치
    'RWRA':  5573,   # 오른쪽 손목 A
    'RWRB':  5568,   # 오른쪽 손목 B
    'LHIP':  1806,   # 왼쪽 고관절
    'RHIP':  6573,   # 오른쪽 고관절
    'RKNEE': 4495,   # 오른쪽 무릎
    'RANK':  6727,   # 오른쪽 발목
    'LKNEE': 1069,   # 왼쪽 무릎
    'LANK':  3327,   # 왼쪽 발목
    # ... 43개 전체
}
\end{verbatim}
\end{codebox}

\section{TRC 파일 생성 + OpenSim IK}

\begin{codebox}[title={SMPL $\to$ TRC $\to$ OpenSim}]
\begin{verbatim}
# 1. SMPL FK → 꼭짓점 추출
out = smpl_model(body_pose=bp, betas=betas, ...)
vertices = out.vertices[0].numpy()  # (6890, 3)

# 2. 마커 위치 추출
markers = {name: vertices[idx] for name, idx in MARKER_MAP.items()}

# 3. TRC 파일 작성
write_trc(markers, fps=240, output='markers.trc')

# 4. OpenSim IK
# 모델: LaiUhlrich2022.osim (SAM4Dcap에서 제공)
# opensim.InverseKinematicsTool → .mot 파일 (관절 각도)
\end{verbatim}
\end{codebox}


% ====================================================================
\chapter{전체 시간표}
% ====================================================================

\begin{center}
\renewcommand{\arraystretch}{1.4}
\begin{tabular}{clcc}
\toprule
\textbf{순서} & \textbf{단계} & \textbf{시간} & \textbf{누적} \\
\midrule
1 & SAM 3D Body 설치 (conda + pip + checkpoint) & 20--30분 & 30분 \\
2 & 테스트 추론 (cam3, 100프레임) & 2분 & 32분 \\
3a & 전체 추론 (7cam $\times$ 999프레임) & 60--120분 & 2.5시간 \\
3b & MHR $\to$ SMPL 변환 & 10--20분 & 3시간 \\
4 & 7-view reproj fitting ($\lambda$=0.0001) & 20--40분 & 3.5시간 \\
5 & GART 7-view 30K스텝 & 30--60분 & 4.5시간 \\
6 & SMPL $\to$ TRC $\to$ OpenSim IK & 30분 & 5시간 \\
7 & Vicon 검증 (동기화 완료) & 30분 & \textbf{5.5시간} \\
\bottomrule
\end{tabular}
\end{center}

\begin{importantnote}
\textbf{병렬화 가능}: Step 3a에서 7개 카메라를 순차 실행하지만, GPU 메모리가 충분하면 2--3개 프로세스 병렬 실행으로 시간 단축 가능.
\end{importantnote}


% ====================================================================
\chapter{품질 게이트 + 리스크}
% ====================================================================

\section{품질 게이트}

\begin{center}
\renewcommand{\arraystretch}{1.3}
\begin{tabular}{clll}
\toprule
\textbf{게이트} & \textbf{위치} & \textbf{통과 기준} & \textbf{실패 시} \\
\midrule
G0 & Step 1 후 & SAM 3D Body 1장 추론 성공 & pip 디버깅 \\
G1 & Step 2 후 & MHR 메시 시각적 확인 (피팅 합리적) & bbox/crop 조정 \\
G2 & Step 3 후 & SMPL 변환 성공, 꼭짓점 수 6,890 & 방법 B 전환 \\
G3 & Step 4 후 & reproj $<$ 100px & $\lambda$ 재조정 \\
G4 & Step 5 후 & PSNR $>$ 20dB & 스텝 증가 \\
G5 & Step 6 후 & OpenSim IK 수렴 & 마커 매핑 재조정 \\
\bottomrule
\end{tabular}
\end{center}

\section{리스크}

\begin{center}
\renewcommand{\arraystretch}{1.3}
\begin{tabular}{p{4cm}ccp{4.5cm}}
\toprule
\textbf{리스크} & \textbf{확률} & \textbf{영향} & \textbf{대응} \\
\midrule
SAM 3D Body detectron2 빌드 실패 & 낮음 & 높음 & PyTorch 버전 맞춤 or 다른 커밋 \\
MHR$\to$SMPL 변환 발산 & 중간 & 중간 & 직접 vertex 피팅 (방법 B) \\
pyrender headless 실패 & 중간 & 낮음 & EGL/osmesa 설정 \\
추론 속도 $<$0.5 FPS & 낮음 & 중간 & cam 수 4개로 축소 \\
cam1 세로에서 검출 실패 & 낮음 & 낮음 & cam1 제외 (6-view) \\
\bottomrule
\end{tabular}
\end{center}


% ====================================================================
\chapter{U-HMR과의 A/B 비교 실험}
% ====================================================================

SAM 3D Body가 정말로 U-HMR보다 나은지 정량적으로 비교합니다.

\section{비교 설계}

\begin{center}
\renewcommand{\arraystretch}{1.3}
\begin{tabular}{lcc}
\toprule
& \textbf{경로 A (현재)} & \textbf{경로 B (SAM 3D Body)} \\
\midrule
초기값 & U-HMR 4-view & SAM 3D Body 7-view 중위수 \\
캘리브레이션 & VGGT & VGGT (reproj용) \\
Reproj fitting & 동일 스크립트, 동일 $\lambda$ & 동일 \\
GART & 동일 설정 & 동일 \\
\midrule
비교 지표 & \multicolumn{2}{c}{Reproj error, MPJPE vs Vicon, PSNR} \\
\bottomrule
\end{tabular}
\end{center}

\begin{keyidea}
\textbf{유일한 변수}: 초기 body\_pose의 출처.\\
나머지 파이프라인을 동일하게 유지하여 SAM 3D Body의 순수 기여도를 측정합니다.\\
cam3 100프레임으로 빠르게 비교 후, 전체 실행 여부를 결정합니다.
\end{keyidea}

\section{빠른 비교 프로토콜 (30분)}

\begin{enumerate}
\item cam3 100프레임에 SAM 3D Body 추론 (2분)
\item MHR $\to$ SMPL 변환 (2분)
\item cam3 단안 reproj 오차 계산 (1분)
\item U-HMR cam3 결과와 비교 (이미 존재)
\item \textbf{SAM 3D Body reproj $<$ U-HMR reproj 이면} $\to$ 전체 실행 진행
\end{enumerate}


% ====================================================================
\chapter{서버 실행 체크리스트}
% ====================================================================

\begin{checklistbox}[title={서버 접속 후 실행 순서}]
\begin{enumerate}
\item[$\square$] \texttt{conda create -n sam3d python=3.11 -y}
\item[$\square$] \texttt{pip install torch ...} (PyTorch + 의존성)
\item[$\square$] \texttt{pip install detectron2} (소스 빌드)
\item[$\square$] HuggingFace 게이트 승인 + 체크포인트 다운로드 (2.8GB)
\item[$\square$] \texttt{git clone sam-3d-body} + 테스트 추론 1장
\item[$\square$] \textbf{G0 게이트}: MHR 메시 (18,439 vtx) 출력 확인
\item[$\square$] cam3 100프레임 추론 $\to$ U-HMR 비교 (A/B 테스트)
\item[$\square$] \textbf{G1 게이트}: SAM 3D Body $\geq$ U-HMR 확인
\item[$\square$] 7-cam 전체 추론 (999프레임)
\item[$\square$] MHR $\to$ SMPL 변환
\item[$\square$] \textbf{G2 게이트}: SMPL 꼭짓점 6,890 확인
\item[$\square$] 7-view reproj fitting ($\lambda$=0.0001)
\item[$\square$] \textbf{G3 게이트}: reproj $<$ 100px
\item[$\square$] GART 7-view 30K스텝 학습
\item[$\square$] \textbf{G4 게이트}: PSNR $>$ 20dB
\item[$\square$] SMPL $\to$ TRC $\to$ OpenSim IK
\item[$\square$] Vicon 검증 (MPJPE, 관절각 RMSE)
\end{enumerate}
\end{checklistbox}

\begin{summary}
\textbf{SAM 3D Body 파이프라인 요약}:
\begin{itemize}
\item 설치 20분, 캘리브레이션 의존성 제거, 7-cam 모두 활용 가능
\item PA-MPJPE 30.4mm (SOTA) 단안 추론 $\times$ 7개 카메라
\item MHR 18,439 꼭짓점 $\to$ SMPL 6,890 변환 $\to$ GART 입력
\item 기존 reproj fitting + GART 스크립트 재활용
\item SAM4Dcap 마커 매핑으로 OpenSim 브릿지
\item Vicon 동기화 이미 완료 (125프레임 오프셋)
\item \textbf{총 5.5시간, 6개 품질 게이트}
\end{itemize}
\end{summary}

\end{document}